% Part: lambda-calculus
% Chapter: church-rosser
% Section: definitions-and-properties

\documentclass[../../../include/open-logic-section]{subfiles}

\begin{document}

\olfileid[id]{lam}{cr}{dap}

\olsection{Definisi dan Sifat}

Dalam bab ini kita memperkenalkan konsep sifat Church--Rosser dan beberapa
sifat umumnya.

\begin{defn}[Sifat Church--Rosser, CR] Suatu relasi $\xredone$ pada suku
  dikatakan memenuhi \emph{sifat Church--Rosser} jika dan hanya jika, setiap
  kali $M \xredone P$ dan $M \xredone Q$, terdapat suatu~$N$ sedemikian
  sehingga $P \xredone N$ dan $Q \xredone N$.
\end{defn}

Kalkulus lambda dapat dipandang sebagai model komputasi yang memperlakukan
suku-suku dalam bentuk normal sebagai ``nilai'' dan relasi ketereduksian pada
suku sebagai ``aturan perhitungan.'' Sifat Church--Rosser menyatakan bahwa
apabila terdapat lebih dari satu cara untuk melanjutkan suatu perhitungan,
ekspresi itu tetap hanya mempunyai satu nilai akhir.

Sebagai contoh dari aljabar elementer, terdapat lebih dari satu cara untuk
menghitung $4 \times (1+2)+3$. Ekspresi ini dapat direduksi menjadi
$4\times3+3$ (jika kita mula-mula mereduksi $1+2$ menjadi~$3$), atau menjadi
$4\times1+4\times2+3$ (jika kita mula-mula menggunakan distributivitas pada
$4\times(1+2)$). Namun, keduanya dapat direduksi lebih lanjut menjadi~$12+3$.

Jika $\xredone$ diambil sebagai reduksi-$\beta$, mudah terlihat bahwa salah
satu konsekuensi sifat Church--Rosser ialah bahwa, jika suatu suku mempunyai
bentuk normal, bentuk itu tunggal. Andaikan $M$ dapat direduksi menjadi $P$
dan $Q$, yang keduanya berbentuk normal. Berdasarkan sifat Church--Rosser,
terdapat suatu $N$ yang menjadi hasil reduksi baik dari $P$ maupun dari $Q$.
Karena menurut asumsi $P$ dan $Q$ berbentuk normal, reduksi dari $P$ dan $Q$
ke $N$ hanya mungkin merupakan reduksi trivial; yaitu, $P$, $Q$, dan $N$
identik. Hal ini membenarkan penggunaan istilah \emph{bentuk normal} suatu
suku.

Dalam pandangan kalkulus lambda sebagai model komputasi, bentuk normal suatu
suku dapat dianggap sebagai ``hasil akhir'' komputasi yang dimulai dari suku
itu. Korolari di atas berarti bahwa suatu komputasi mempunyai paling banyak
satu hasil akhir, sama seperti perhitungan $4\times(1+2)+3$ hanya mempunyai
satu hasil, yaitu~$15$.

\begin{thm} \ollabel{thm:str}
  Jika suatu relasi $\xredone$ memenuhi sifat Church--Rosser dan $\xred$
  adalah relasi transitif terkecil yang memuat $\xredone$, maka $\xred$ juga
  memenuhi sifat Church--Rosser.
\end{thm}

\begin{proof}
  Andaikan
  \begin{align*}
    M & \xredone P_1 \xredone \dots \xredone P_m \text{ dan}\\
    M & \xredone Q_1 \xredone \dots \xredone Q_n.
  \end{align*}
  Kita membuktikan teorema dengan membangun kisi $N$ yang terdiri atas suku,
  dengan tinggi $m+1$ dan lebar $n+1$. Kita menggunakan $N_{i,j}$ untuk
  menyatakan suku pada baris ke-$i$ dan kolom ke-$j$.

  Kita membangun $N$ sedemikian sehingga $N_{i,j}\xredone N_{i+1,j}$ dan
  $N_{i,j}\xredone N_{i,j+1}$. Kisi itu didefinisikan sebagai berikut:
  \begin{align*}
    N_{0,0} &= M \\
    N_{i,0} &= P_i && \text{jika } 1 \le i \le m \\
    N_{0,j} &= Q_j && \text{jika } 1 \le j \le n
  \intertext{dan dalam kasus lainnya:}
    N_{i,j} &= R,
  \end{align*}
  dengan $R$ suatu suku sedemikian sehingga $N_{i-1,j}\xredone R$ dan
  $N_{i,j-1}\xredone R$. Berdasarkan sifat Church--Rosser bagi $\xredone$,
  suku semacam itu selalu ada.

  Sekarang kita mempunyai $N_{m,0}\xredone\dots\xredone N_{m,n}$ dan
  $N_{0,n}\xredone\dots\xredone N_{m,n}$. Perhatikan bahwa $N_{m,0}=P$
  dan $N_{0,n}=Q$. Teorema mengikuti dari definisi~$\xred$.
\end{proof}

\end{document}
